\documentclass{article}
\usepackage[T1]{fontenc}
\usepackage{lmodern}
\usepackage{graphicx}
\usepackage{amsmath,amssymb}
\usepackage[a4paper,margin=2cm]{geometry}
\usepackage[absolute,overlay]{textpos}
\setlength{\TPHorizModule}{1mm}
\setlength{\TPVertModule}{1mm}
\usepackage[indonesian]{babel}
\usepackage{hyperref}
\setlength{\emergencystretch}{2em}
% unit: Halaman 2

\begin{document}

% === Halaman 2 (python/native) ===

2

Pendahuluan

• RSA merupakan algoritma kunci-publik yang paling terkenal dan 

paling banyak aplikasinya.

• Dibuat oleh tiga peneliti dari MIT (Massachussets Institute of 

Technology), yaitu Ronald Rivest, Adi Shamir, dan Leonard Adleman, 

pada tahun 1976. 

• RSA = Rivest-Shamir-Adleman

• Keamanan algoritma RSA terletak pada sulitnya memfaktorkan 

bilangan bulat yang besar menjadi faktor-faktor prima.

\end{document}
